\section{\texttt{try}-with-resources}
\label{sec:twr}

\lstinputlisting[language=Java]{generated/source/TryWithResources.java}

\inputbytecode{TryWithResources}

\paragraph{Lendo a listagem.}
\begin{itemize}
\item Offsets 26--30: caminho normal, \texttt{close()} (27) e salto para
  o fim (\texttt{goto 49}, offset 30).
\item A \texttt{Exception table} diz: exceção nos offsets 8 a 25 (o
  \texttt{to}~26 é exclusivo, e o \texttt{close()} em 27 fica fora) vai
  para 33. Ali a exceção principal é guardada (\texttt{astore\_1}) e
  \texttt{close()} é chamado de novo.
\item Se esse \texttt{close()} também lança, a segunda linha da tabela
  leva ao offset 41, que chama
  \texttt{Throwable.addSuppressed}: a exceção do \texttt{close()} é
  pendurada na principal, que é relançada (\texttt{athrow}, offset 48).
\end{itemize}

É a tradução que a JLS (§14.20.3.1) especifica, com a variável
\texttt{\#primaryExc} e \texttt{addSuppressed}. Compare com
\texttt{ThreadsVirtuais} (seção~\ref{sec:virtual-threads}): lá o recurso
vem de uma chamada de método e o \texttt{javac} gerou \texttt{ifnull}
antes de cada \texttt{close()}, como a JLS pede; aqui o recurso vem de
um \texttt{new}, que nunca é \texttt{null}, e o teste não aparece.

\paragraph{O bug que isso explica.}
Com \texttt{try}/\texttt{finally} escrito à mão, uma exceção no
\texttt{close()} substitui a exceção original, e o log mostra
\enquote{erro ao fechar conexão} em vez da causa real. O
\texttt{try}-with-resources preserva a causa e anexa a outra como
\emph{suppressed}: procure \texttt{Suppressed:} no stack trace. É o
argumento do item 9 do \citetitle{livro:bloch2017} (\enquote{Prefer
try-with-resources to try-finally}).
